\documentclass[11pt]{article}

\usepackage[a4paper,margin=2.3cm]{geometry}
\usepackage[utf8]{inputenc}
\usepackage[T1]{fontenc}
\usepackage{booktabs}
\usepackage{array}
\usepackage{longtable}
\usepackage{xcolor}
\usepackage{hyperref}
\usepackage{parskip}
\usepackage{enumitem}
\usepackage{amsmath}
\usepackage{listings}
\usepackage{caption}

\hypersetup{colorlinks=true, linkcolor=blue, urlcolor=blue, citecolor=blue}

\definecolor{codebg}{gray}{0.95}
\lstset{
  basicstyle=\ttfamily\small,
  backgroundcolor=\color{codebg},
  breaklines=true,
  frame=single,
  columns=fullflexible
}

\title{\textbf{TIES-Merging Notebooks: Comparative Analysis}\\
\large ties-merging-fast-1 \,$\to$\, ties-merging-fast-4}
\author{}
\date{}

\begin{document}
\maketitle
\vspace{-1.5em}

\section{Overview}

The four files represent successive iterations of the \emph{same} pipeline,
each one responding to a problem diagnosed in the previous run. The
UNet/DDPM architecture code (cell~1 in notebooks 1--3, and the
\texttt{architecture.py} module in notebook 4) is byte-for-byte identical
across all four notebooks — only the configuration, merge logic, evaluation
protocol, and search strategy change.

\section{Notebook-by-notebook summary}

\subsection*{1. \texttt{ties-merging-fast-1.ipynb} --- baseline pipeline}
The first, ``clean'' version of the pipeline. Seven cells: architecture,
configuration, merge functions (\texttt{trim\_global}, sign election, TIES
merge with \texttt{mean} aggregation only), evaluation functions, a two-stage
grid search (proxy $v$-loss ranking $\to$ PSNR/SSIM + retention on the
top-$K$), final evaluation, and visualization. The grid search only sweeps
\texttt{keep\_ratio} $\in\{0.1,0.2,0.3,0.5\}$, \texttt{lambda\_scale}
$\in\{0.4,0.6,0.8,1.0,1.3\}$ and three \texttt{task\_lambdas} settings, all
combined with the classical TIES \texttt{mean} rule. Grid search and
evaluation both run on the validation splits of the three training
manifests.

\subsection*{2. \texttt{ties-merging-fast-2.ipynb} --- diagnostic pass}
Structurally identical to notebook 1 (same 7 cells, same numbers wherever
they overlap), but with one extra cell inserted after the evaluation
functions: a \textbf{diagnostic block} that stress-tests four hypotheses for
why the merge underperforms:
\begin{itemize}[nosep]
  \item \textbf{H1} -- the merge itself is the problem (task-vector interference);
  \item \textbf{H2} -- the specialists are undertrained;
  \item \textbf{H3} -- there is a code/sampler bug;
  \item \textbf{H4} -- the evaluation metric is uninformative.
\end{itemize}
It runs a trivial-baseline test (identity mapping vs.\ target), and a
DDIM-vs-native-sampler comparison for each specialist. The conclusion,
visible in the printed output, is that \textbf{all three specialists work
correctly} (e.g.\ star removal reaches 40.6\,dB with either sampler, far
above the 23.4\,dB identity baseline) — which rules out H2/H3/H4 and isolates
the problem to the merge step itself (H1). This diagnosis is what motivates
notebook 3.

\subsection*{3. \texttt{ties-merging-fast-3.ipynb} --- ``v5'', addressing H1}
Rewritten based on the notebook-2 diagnosis. Key changes:
\begin{itemize}[nosep]
  \item A new \texttt{merge\_method} knob: besides classical \texttt{mean}
  (which divides the merged contribution by the number of agreeing tasks and
  was identified as diluting the signal), it adds \texttt{sum} (preserves
  task-vector magnitude) and \texttt{dropout\_rescale} (DARE-style trim +
  rescale by $1/\text{keep\_ratio}$).
  \item A stronger grid: \texttt{keep\_ratio} up to 0.8, \texttt{lambda\_scale}
  up to 1.6, because a weak $\theta_0$ base (only $\sim$12--23\,dB before
  merging) needs larger corrections.
  \item Selection criterion changed to mean retention with \emph{min}-retention
  as a tie-break, to avoid rewarding configurations that sacrifice one task
  entirely.
  \item All \texttt{dataset\_merged}-related code removed; evaluation is
  exclusively on the three manifest validation splits (same protocol as 1/2,
  narrower scope).
\end{itemize}
In practice only \texttt{merge\_method="sum"} was actually swept (the other
two options are present in code but commented out of the grid), over
\texttt{keep\_ratio}$\in\{0.2,0.5,0.8\}$ and
\texttt{lambda\_scale}$\in\{0.5,0.8\}$ (15 configurations total). The
notebook completed end-to-end and saved \texttt{ties.pth}.

\subsection*{4. \texttt{ties-merging-fast-4.ipynb} --- ``v6'', task-arithmetic rewrite, multi-GPU}
A structural rewrite rather than an incremental patch. Instead of one
monolithic notebook it \emph{generates a Python package}
(\texttt{/kaggle/working/ties\_merge/\{architecture,configuration,merging,\allowbreak
evaluation,search\}.py}) via \texttt{\%\%writefile} cells, and launches the
search as a separate multi-GPU process with
\texttt{accelerate launch --num\_processes=2}. Design changes:
\begin{itemize}[nosep]
  \item \textbf{Dataset}: abandons the training manifests in favor of the
  \texttt{dataset\_merged/\{SU,IR,SR,\dots\}} folders (same data
  \texttt{task-arithmetic} uses), with a \texttt{subsample\_pct} (10\,\%) for
  speed, read through a new \texttt{DegradedImageDataset}.
  \item \textbf{Two-stage search with a different metric}: Stage 1 evaluates
  every \texttt{(merge\_method, task\_lambdas, keep\_ratio, lambda\_scale)}
  combination on the three \emph{single-task} folders using a
  normalized score $\text{score}=\text{mean}_t(\text{MSE}_{\text{merge}, t} /
  \text{MSE}_{\text{expert}, t})$, discarding any config whose per-task ratio
  exceeds \texttt{stage1\_max\_task\_ratio = 2.5} (a ``safe zone'' filter).
  Stage 2 was intended to re-rank the surviving ``safe'' configs on
  \emph{multi-degradation} folders (\texttt{IR\_SU}, \texttt{SR\_SU},
  \texttt{SR\_IR}, \texttt{SR\_IR\_SU}) by mean MSE, and save the winner.
  \item TIES machinery itself (task vector, global trim, sign election,
  \texttt{merge\_method}) is unchanged from notebook 3; only the search
  protocol and dataset source come from \texttt{task-arithmetic}.
\end{itemize}
\textbf{Execution outcome:} Stage 1 completed successfully (15/15
configurations evaluated, 2 passed the safe-zone filter). Stage 2 crashed
immediately with
% \lstinline{FileNotFoundError: .../dataset_merged/IR_SR} — the multi-task
% folder the code tried to read (\texttt{IR\_SR}) does not match the naming
% used in \texttt{CONFIG["tasks\_stage2"]} (\texttt{SR\_IR}), i.e.\ a
% path/naming-convention bug in Stage~2's folder resolution. As a result the
% \texttt{accelerate} job terminated with a non-zero exit code, \textbf{no
% \texttt{ties.pth} was ever written}, and the final in-notebook visualization
% cell (cell~8) fails with \lstinline{FileNotFoundError:
% '/kaggle/working/ties.pth'}.

\section{Cross-cutting differences}

\begin{itemize}[nosep]
  \item \textbf{Architecture}: identical in all four (UNet + DDPM,
  \texttt{v}-prediction, 200 diffusion steps, EMA weights preferred).
  \item \textbf{Merge core (trim + sign election)}: unchanged since notebook 1;
  what changes is how the per-task, sign-agreed deltas are \emph{combined}
  (\texttt{mean} only in 1/2, \texttt{mean}/\texttt{sum}/\texttt{dropout\_rescale}
  available from 3 onward, though only \texttt{sum} is swept in 3 and 4).
  \item \textbf{Evaluation data}: notebooks 1--3 use the \emph{training
  manifests'} validation splits (same distribution as training, in-distribution
  generalization check). Notebook 4 switches to the held-out
  \texttt{dataset\_merged} folders, including genuinely novel multi-degradation
  combinations — closer to the original stated goal of testing composition of
  unseen degradations, but this is also where it fails.
  \item \textbf{Search protocol}: 1/2/3 use a 2-stage proxy-$v$-loss $\to$
  PSNR/SSIM-retention search over the \emph{same} single-task validation data.
  4 uses a differently-motivated 2-stage search (single-task safe-zone filter
  $\to$ multi-task re-ranking) over a \emph{different} dataset per stage.
  \item \textbf{Selection metric}: retention relative to the specialist
  (1, 2, 3) vs. MSE ratio relative to the expert, unnormalized by $\theta_0$
  (4).
  \item \textbf{Infrastructure}: 1--3 are single-process, in-notebook. 4 is a
  generated multi-file package executed via \texttt{accelerate} on 2 GPUs,
  which is also the source of its failure (Stage-2 dataset path only checked
  at run time, in a separate process, uncaught by any prior validation cell).
  \item \textbf{Deliverable}: 1, 2, and 3 each produce a working
  \texttt{ties.pth} and a completed visualization cell. 4 does not produce a
  merged checkpoint in the run captured in this notebook, because Stage 2
  never completes.
\end{itemize}

\section{Summary table}

\begin{table}[htbp]
\centering
\small
\renewcommand{\arraystretch}{1.3}
\begin{tabular}{@{}p{1.7cm}p{4.3cm}p{4.3cm}p{4.9cm}@{}}
\toprule
\textbf{Notebook} & \textbf{Key design change vs.\ previous} & \textbf{Search / metric} & \textbf{Captured output} \\
\midrule

\textbf{1} \newline (baseline) &
Reference TIES pipeline. Only \texttt{mean} merge rule. &
2-stage: proxy $v$-loss $\to$ PSNR/SSIM + retention, on manifest val splits. &
Ran end-to-end. $\theta_0$: 23--25\,dB; specialists 26--41\,dB. Merged model: restoration 28.99\,dB, super-res 23.50\,dB, star removal 21.82\,dB (SSIM 0.55--0.88) — \textbf{below $\theta_0$ on every task} (deltas $-0.97$ to $-3.63$\,dB). \\

\textbf{2} \newline (diagnosis) &
Same pipeline as 1 (identical numbers), plus an inserted \textbf{diagnostic cell} (H1--H4 hypothesis tests). &
Same as 1, plus trivial-baseline \& DDIM-vs-native sampler tests. &
Identity baselines 23--25\,dB; specialists 26--41\,dB with \emph{either} sampler $\Rightarrow$ \textbf{specialists confirmed healthy}; problem isolated to the merge (H1). Final merge numbers identical to notebook 1. \\

\textbf{3} \newline (``v5'') &
New \texttt{merge\_method} (\texttt{sum}/\texttt{dropout\_rescale} added to \texttt{mean}); wider grid (\texttt{keep\_ratio}$\le$0.8, $\lambda\le$1.6); \texttt{dataset\_merged} code removed. &
Same 2-stage protocol as 1/2, but only \texttt{sum} swept; selection = mean retention, min-retention tie-break. &
Ran end-to-end (15 configs, $\sim$315\,s). Best: \texttt{sum}, keep=0.8, $\lambda$=0.5 $\to$ PSNR restoration 31.6\,dB, super-res 23.5\,dB, star removal 24.9\,dB; retention mean $-0.18$, min $-1.09$ (still negative overall, but a clear improvement over notebook 1's merge, especially for restoration). \\

\textbf{4} \newline (``v6'', best design) &
Full rewrite as a Python package + \texttt{accelerate}, multi-GPU. Dataset switched to \texttt{dataset\_merged} folders (incl.\ true unseen multi-degradation combos, e.g.\ \texttt{SR\_IR\_SU}); normalized MSE-ratio scoring; safe-zone filter (ratio $\le$2.5) before Stage 2. &
Stage 1: MSE ratio vs.\ expert on single tasks, 15 configs. Stage 2 (intended): mean MSE on 4 multi-task folders. &
Stage 1 completed: best safe configs at keep=0.5/0.8, $\lambda$=0.5 (score 1.00 / 1.20); expert baselines 21--37\,dB. \textbf{Stage 2 crashed} (\texttt{FileNotFoundError} on folder \texttt{IR\_SR} vs.\ configured \texttt{SR\_IR}) $\Rightarrow$ no \texttt{ties.pth} saved; visualization cell fails. Design is the most complete/correct in principle, but \textbf{this run did not finish}. \\

\bottomrule
\end{tabular}
\caption{Comparison of the four TIES-merging notebook iterations: design changes, search/evaluation protocol, and the numbers actually produced by the captured run of each notebook.}
\end{table}

\section{Note on ``notebook 4 is the best''}

By \textbf{design}, notebook 4 is the most advanced iteration: it evaluates on
genuinely held-out data (including multi-degradation combinations never seen
by any single specialist), uses a cleaner normalized scoring function, adds a
safe-zone filter against destructive merges, and is engineered to scale to
multi-GPU search. That said, the run captured in this file did \emph{not}
complete: Stage~2 fails on a folder-naming mismatch
(\texttt{tasks\_stage2} configured as \texttt{SR\_IR} vs.\ an on-disk/expected
folder \texttt{IR\_SR}) before any multi-task score is computed, so no merged
checkpoint (\texttt{ties.pth}) exists at the end of the notebook and the final
visualization cell errors out. Fixing that one path (or renaming the folder to
match) is a small change, and would let Stage~2 -- and therefore the
qualitatively strongest version of the pipeline -- actually run to completion.
Notebook 3 is, among the four, the only later iteration with both an improved
merge strategy \emph{and} a completed, saved output, making it the best
\emph{currently working} checkpoint, while notebook 4 is the best \emph{design}
once its Stage-2 path bug is fixed.

\end{document}
